\subsection{Two lower roots at a finite upper saddle}
\label{rw:finite:three-saddle}

When both simple lower roots meet an upper saddle, all three
passages can lose their exits together. Fix $\bar B=-\beta$,
$0<\beta_*\le\beta\le\beta^*$, near the triangle center
$\bar m=\beta$, $\bar c=-\bar a$, $\nu=0$. The actual upper
ratio $r$ stays in a fixed compact subset of $(0,\infty)$. Retaining
both lower deficits as well as this ratio allows all three passages
to degenerate independently. Use the actual splitting coordinates of Lemma~\ref{rw:coordinates:finite-splittings}, and the
upper coordinates and regular legs of
Lemma~\ref{rw:coordinates:two-center-upper} near the reversible
intersection. On the remaining compact parts, extend just the triangle
coordinates $L_\pm L_0^{\epsilon\beta}$ from
Lemma~\ref{rw:centers:triangle}, and straighten and normalize the
neighboring axes by the same construction. The axis graph denominators
have positive margins on every compact positive-ratio interval; no
resonant shear is needed for this extension. Its orbital remainder is
$\epsilon$ times a bounded function and vanishes on the triangle
component. The same regular outer legs therefore give
\[
 u=G_+(t)=u_0+tA_0(t),\qquad
 z=G_-(v)=z_0+vB_0(v),\qquad v=t^re^{w(t)},
\]
where $A_0,B_0,G_+',G_-'>c>0$. The lower maps are those of
Proposition~\ref{rw:lower-preparation}:
\begin{equation}\label{rw:finite:three-powers}
 D_\pm(s)=K_\pm s^{h_\pm/\epsilon}e^{j_\pm(s)},\qquad
 K_\pm>0,\quad h_\pm>h_*>0,\quad j_\pm(0)=0.
\end{equation}
The upper passage and both lower deficit passages define the signed
central comparison on
\begin{equation}\label{rw:finite:three-domain}
 0<t<t_0,\qquad u>0,\qquad z>0.
\end{equation}
Since $u,z$ increase, this is an interval. The closing function there is
\[
 F(t)=D_+(u)-D_-(z)-\sqrt\epsilon\,\nu H,\qquad H>h_*>0.
\]
The weak references, and hence $H$, are independent of $t$.
The complete central outputs are $s_++D_+(u)$ and $s_-+D_-(z)$.
Their positivity is an additional requirement for an actual return,
beyond the deficit conditions in \eqref{rw:finite:three-domain}.
The Gaussian reference calculation in
Lemma~\ref{rw:finite:regular-lower} applies on the fixed inner cores:
$F_\pm$ have positive margins there even when the outer roots are
lost. It gives $H=\sqrt{2\pi}+O(\sqrt\epsilon)$ and, more precisely,
\[
 s_+=-\sqrt\epsilon\,\nu H_+,
 \qquad s_-=\sqrt\epsilon\,\nu H_-,
 \qquad H_+,H_->c>0,\qquad H=H_++H_-.
\]
At a zero of $F$ the two outputs agree. For $\nu\ge0$ the
incoming one is positive, and for $\nu\le0$ the outgoing one is
positive, since $D_\pm>0$. Thus their common output is always
positive at a closing zero. In each core the original height
$\mathfrak z=1+y$ has derivative of fixed sign on $\mathfrak z=0$,
namely the sign of $\nu$ (with the reflected sign on the other
core). An orbit with positive entrance and exit cannot cross that
line and return. Together with the strict deficit passages and finite
transverse connections, this proves that each closing zero gives an actual
positive-height return. At nonclosing inputs we continue to use the
signed comparison.

The numerator calculation below uses coefficient division as well as
input differentiation. Reserve twelve total finite orders in the lower
and upper preparations before shrinking their construction domains.
This supplies six ordinary input and two coefficient derivatives of the regular connectors and
lower relative coefficients, four Euler input and four coefficient
derivatives of the upper relative logarithm, and the two coefficient
derivatives spent in product division below. The lower graph can be
constructed at order $23$ by the $N+11$ reserve in
Proposition~\ref{rw:lower-preparation}. Scale derivatives are Euler
derivatives. The ordinary input derivatives of the lower connectors
are not asserted for the singular composite $t^r$ at zero.

Use the common signed relative extension and absolute zero extension
after Proposition~\ref{rw:generator-division} when a coefficient
segment changes a deficit's sign.
For any fixed desired $S_*>0$, first reduce the input radius $d<1$
so that a half of $d^{h_*/\epsilon}$ absorbs the upper bound for
$\exp(\epsilon\log K_\pm/\epsilon)$, and then reduce $\epsilon$
so that the remaining half-power absorbs all reserved derivatives.
This gives
\begin{equation}\label{rw:finite:three-flat}
 |\partial_s^i\partial_{\mathbf b}^{\alpha}\widehat D_\pm|
       \le C\epsilon^{-M}e^{-S_*/\epsilon}
\end{equation}
through the required finite orders. These extensions preserve both
center identities by that construction.
Keep the original $t$ fixed on the coefficient segments and impose
central positivity on actual returns.

Put
\[
 \mathcal K=\epsilon\log
 \frac{D_-'(z)G_-'(v)v'}{D_+'(u)G_+'(t)},\qquad
 D=t\partial_t,\qquad \kappa=r-1,\qquad
 E=\frac{tv''}{v'}-\kappa.
\]
Ordinary primes on $G_\pm$ and $D_\pm$ mean their own input
derivatives. For use in the chain rule define
\[
 R_\pm(s)=j_\pm'(s)+
       \frac{j_\pm'(s)+sj_\pm''(s)}{h_\pm/\epsilon+sj_\pm'(s)}.
\]
Then $D_\pm''/D_\pm'=(h_\pm/\epsilon-1)/s+R_\pm$ on $s>0$.
The displayed formula also defines bounded auxiliary coefficients on
the signed interval; no logarithm of a zero extension is taken.

\begin{proposition}[Complete three-saddle numerator]
\label{rw:finite:three-numerator}
On the comparison interval the following function is exactly
$tuz\,\mathcal K'$:
\begin{equation}\label{rw:finite:three-M}
\begin{aligned}
 M={}&(h_--\epsilon)uDz-(h_+-\epsilon)zDu
              +\epsilon uz(\kappa+E)\\
 &+\epsilon uz\left[
 R_-(z)Dz-R_+(u)Du+
 \frac{G_-''(v)}{G_-'(v)}Dv-
 \frac{G_+''(t)}{G_+'(t)}t\right].
\end{aligned}
\end{equation}
This formula defines its auxiliary extension. Its endpoint value is
\[
 M(0+)=\epsilon\kappa u_0z_0.
\]
Near ratio one put
$\xi_1=(u_0+z_0)/2$, $\xi_2=(u_0-z_0)/2$ and
$I=|\xi_2|+|\kappa\xi_1|$. Then
\begin{equation}\label{rw:finite:three-comp}
\begin{aligned}
 M&=C_0+Tt-Wt\omega(t;\kappa)+\mathcal R_3,\\
 C_0&=\epsilon\kappa u_0z_0,\\
 P&=u_0z_1(rh_--\epsilon)+\epsilon u_0z_0rJ_-,\\
 N&=z_0u_1(h_+-\epsilon r)+\epsilon u_0z_0J_+,\\
 T&=P-N+\epsilon u_0z_0\sigma_s(2+1/r),\\
 W&=\kappa P+2\epsilon u_0z_0\sigma_s r,
\end{aligned}
\end{equation}
where $u_1=G_+'(0)$, $z_1=G_-'(0)$ and
\[
 J_+=u_1R_+(u_0)+G_+''(0)/u_1,\qquad
 J_-=z_1R_-(z_0)+G_-''(0)/z_1.
\]
Here $\sigma_s$ denotes the resonant coefficient $\sigma$ of
\eqref{rw:upper:comp-map} in the present center-preserving coordinates.
For $|\kappa|\le1/64$,
\begin{equation}\label{rw:finite:three-remainder}
 |D^j\mathcal R_3|\le C(I+|\nu|)t^{7/4},\qquad 0\le j\le2.
\end{equation}
If this fibre has one actual closing zero, then
\begin{equation}\label{rw:finite:three-forcing}
 |\nu|\le\delta_\epsilon I,\qquad
 \delta_\epsilon=C\epsilon^{-M-1/2}e^{-S_*/\epsilon}\longrightarrow0.
\end{equation}
\end{proposition}

\begin{proof}
Differentiate $\mathcal K$ before clearing its denominators.
The two power rates give the first two terms in
\eqref{rw:finite:three-M}; the upper second derivative gives
$\epsilon uz(\kappa+E)$. The ordinary curvature of each regular leg
gives the last two terms in its bracket. Thus the numerator retains
every coordinate contribution to the complete central derivative.

To verify the expansion, first regard $t$ and $t^r$ as independent
small quantities. The endpoint Taylor formulas give
\[
 u=u_0+u_1t+O(t^2),\qquad
 z=z_0+z_1t^r+\text{terms of order greater than one}.
\]
The terms linear in these two quantities are
\[
 C_0+Pt^r-Nt-\epsilon u_0z_0\sigma_s Q(t),
 \qquad Q=2rt\omega-(2+1/r)t.
\]
In particular the $J_\pm$ contributions are present even when the
connectors are nonlinear. The exact equality
$t^r=t-\kappa t\omega$ gives \eqref{rw:finite:three-comp}.

We first estimate the remainder uniformly on the auxiliary box,
before dividing by the center generators. Write $\mathcal F_a$ for a
function whose selected
$D^j\partial_{\mathbf b}^{\alpha}$ derivatives are $O(t^a)$,
$j\le2$, $|\alpha|\le2$. With $a=15/16$ the functions
$t,t^r,t\omega$ belong to $\mathcal F_a$ on
$|\kappa|\le1/64$; every fixed parameter logarithm is absorbed by
$63/64-a>0$. The upper preparation gives
$v-t^r,Dv-rt^r\in\mathcal F_{2a}$ and
$E+\sigma_sQ\in\mathcal F_{7/4}$.
Integral Taylor remainders of $G_\pm$, $R_\pm$ and their regular
curvatures then show that every omitted product has at least two
$\mathcal F_a$ factors or the displayed curvature remainder.
As $2a=15/8>7/4$, it follows that
$D^j\partial_{\mathbf b}^{\alpha}\mathcal R_3=O(t^{7/4})$ on
the entire common auxiliary box.

The triangle center gives vanishing of $M$ and its retained
coefficients on $\xi_1=\xi_2=\nu=0$; the reflected extension
identities give the same vanishing on $\xi_2=\kappa=\nu=0$.
These two axes explain the generator pair $(\xi_2,\kappa\xi_1)$.
Apply the product-ideal division of
Proposition~\ref{rw:generator-division} to $\mathcal R_3$.
Its explicit mixed integral in $\xi_1,\kappa$ spends two coefficient
derivatives and preserves $t^{7/4}$, proving
\eqref{rw:finite:three-remainder}.

Apply the identical division to
$\widehat D_+(G_+(t))-\widehat D_-(G_-(v(t)))$.
Equation~\eqref{rw:finite:three-flat} bounds its divided coefficients
by $C\epsilon^{-M}e^{-S_*/\epsilon}$. At a physical closing zero
Lemma~\ref{rw:centers:relative-forcing} now proves
\eqref{rw:finite:three-forcing}. This is a coefficient constraint
on the whole fibre. If $I=0$, it forces $\nu=0$ and the exact center
identity contributes no isolated zero.
\end{proof}

\paragraph{The scalar compensator calculation.}
The same scalar calculation controls the derivative at a single upper
saddle and the numerator here. Keeping $t\omega$ intact makes its
bounds uniform even when $|\kappa\log t|$ is unbounded.

\begin{lemma}\label{rw:finite:compensator-calculus}
For $|\kappa|\le1/64$ set $r=1+\kappa$ and
\[
 M=C+Tt-Wt\omega(t;\kappa)+R,\qquad L=|T|+|W|>0.
\]
Assume $|R'|\le C_*Lt^{1/2}$ and
$|R''|\le C_*Lt^{-1/2}$. On a common sufficiently small
$(0,t_0)$, $M$ has at most two zeros, with multiplicity.
If it has two, then
\begin{equation}\label{rw:finite:compensator-signs}
 TW>0,\qquad (C+R(0))W>0.
\end{equation}
There is also the following one-zero version. For fixed positive
$c_*,C_*$, if
\[
 W/\kappa\ge c_*,\quad |T|\le C_*|\kappa|,\quad
 |R'|\le C_*|\kappa|t^{1/2},\quad \kappa\ne0,
\]
then first choosing $t_0$ small and then $|\kappa|$ small makes
$M'$ have sign $-\operatorname{sign}\kappa$ throughout the interval.
The constant $C$ is unrestricted in both assertions.
\end{lemma}

\begin{proof}
Let $\phi=(t\omega)'$. Direct differentiation gives
$\phi'=-rt^{\kappa-1}$. On $0<t<e^{-6}$,
$1\le\phi\le16t^{-1/8}$. At any zero of $M'$,
\[
 T=W\phi-R',\qquad
 |W|\ge\frac{L}{2(1+\phi)}
\]
once $C_*t_0^{1/2}<1/2$. Thus $W\ne0$ at such a zero and
\[
 \frac{|R''|}{|W|rt^{\kappa-1}}
 \le \frac{2C_*}{r}(1+\phi)t^{1/2-\kappa}
 \le C C_*t^{23/64}<1/2
\]
after a second common reduction. Every zero of $M'$ is therefore
simple and crosses with sign $W$. There can be at most one such
crossing; Rolle's theorem proves the two-zero bound, including
multiplicity. If $W=0$, the estimate on $R'$ excludes a zero of $M'$.

When $TW\le0$, $|T-W\phi|\ge L>|R'|$, so two zeros are impossible.
If $W>0$, the unique possible critical point is a minimum. Two zeros
force strict decrease to the first of them and hence $M(0+)>0$.
For $W<0$ all signs reverse. The integrable bound on $R'$ supplies
$R(0)$ and proves \eqref{rw:finite:compensator-signs}.
The last assertion is the $\phi$ estimate in the proof of
Lemma~\ref{rw:finite:saddle-counting}\textup{(ii)}, applied to
$M'=T-W\phi+R'$.
\end{proof}

\paragraph{The coefficient columns near ratio one.}
To apply the scalar lemma, we must control its coefficient norm
$|T|+|W|$ by the center generators. Work on the same inverse box
as above. Transversality of $T$ to the reversible component and an
exact factorization of $W$ give the required determinant margin.
At the reversible intersection, let $k_*=G_+'(0)=G_-'(0)>0$.
There $h_+=h_-=1/(2\beta)$ exactly. Put
\[
 A_c=k_*(1/(2\beta)-\epsilon)>c>0.
\]
The coefficient $T$ has $\partial_{\xi_2}T=2A_c$ there.
Moreover the functions $\tau_1$, $\sigma_1$ and $\sigma_s/\epsilon$
vanish on both center axes, so their divisions have bounded columns:
\begin{equation}\label{rw:finite:three-directed-generators}
\begin{aligned}
 \tau_1&=a_\tau\xi_2+b_\tau\kappa\xi_1+c_\tau\nu,
                         &a_\tau&>c>0,\\
 \sigma_s/\epsilon&=-a_s\xi_2+b_s\kappa\xi_1+c_s\nu,
                         &a_s&>c>0,\\
 |\tau_1|+|\sigma_1|&\le C(I+|\nu|).
\end{aligned}
\end{equation}
The positive $a_\tau$ is the inverse of the positive transverse
column in \eqref{rw:coordinates:finite-jacobian}, continued with
its coefficient derivative margin. We verify the second direction
from the original upper field.

The coordinates of Lemma~\ref{rw:coordinates:two-center-upper}
are affine at zero scale, so the first-order calculation in
Appendix~\ref{rw:app:upper-resonant-sign} applies. Absorbing their
positive linear section factors in $\chi_{\rm loc}^2$, it gives on
$\kappa=\nu=0$
\[
 \sigma_s=-\chi_{\rm loc}^2\epsilon\bar a
                  [\bar m^{5/2}+O(\epsilon)],\qquad\chi_{\rm loc}>c.
\]
The factors $\epsilon$ and $\bar a$ are exact, by the affine
zero-scale field and the reversible identity. The finite coordinate
Jacobian gives $\partial_{\bar a}\xi_2>c$ at fixed $\kappa$.
Hence the $\xi_2$ column is strictly negative at the intersection.
The integral division, with the uniform second coefficient jets,
retains that sign on a smaller box and proves
\eqref{rw:finite:three-directed-generators}.

The $W$ column has an exact positive factor on the finite-scale
coefficient box. Define
\[
 Q_*=z_1(rh_--\epsilon)+\epsilon z_0rJ_->c>0,\qquad
 Q^*=u_1(h_+-\epsilon r)+\epsilon u_0J_+>c>0.
\]
Exactly, $P=u_0Q_*$ and $N=z_0Q^*$. On $\xi_2=\nu=0$ write
$\sigma_s=\epsilon b_s\kappa\xi_1$. Then
\begin{equation}\label{rw:finite:three-W-factor}
 W=\kappa\xi_1[Q_*+2\epsilon^2rb_s\xi_1^2].
\end{equation}
The bracket has a uniform positive margin.
Remove $\xi_2,\nu$ by single integral division. On $\nu=0$ this gives
\[
 \binom TW=
 \begin{pmatrix}a_T&b_T\\a_W&b_W\end{pmatrix}
 \binom{\xi_2}{\kappa\xi_1},
 \qquad a_T>c,\quad b_W>c,\quad |b_T|\le C.
\]
The exact formula for $W$ also gives
$|a_W|\le C|\kappa|+C\epsilon^2(|\xi_1|+|\xi_2|)^2$.
Consequently the determinant has a positive margin after reducing
the box; restoring $\nu$ and using
\eqref{rw:finite:three-forcing} yields
\begin{equation}\label{rw:finite:three-column-norm}
 cI\le |T|+|W|\le CI
\end{equation}
on every noncenter fibre with a closing zero. The division uses the
coordinates $\kappa,\xi_1$ themselves; their product
$\kappa\xi_1$ is a generator.

\paragraph{Critical points on the deficit domain.}
The deficit inequalities now constrain a critical point of $M$.
This step uses $u,z>0$ and therefore applies on the comparison
interval, rather than on the whole auxiliary box.

\begin{lemma}\label{rw:finite:three-critical-gaps}
Near ratio one, under \eqref{rw:finite:three-forcing} on a noncenter
fibre, every zero of $M'$ in the comparison interval has $u_0z_0>0$.
\end{lemma}

\begin{proof}
By \eqref{rw:finite:three-comp} and
$1-\kappa(t\omega)'=rt^\kappa$,
\[
 M'=rP t^\kappa-N-\epsilon u_0z_0\sigma_sQ'
                                 +\mathcal R_3'.
\]
Suppose $u_0z_0\le0$, and put
$S(t)=|u_0|t^\kappa+|z_0|$. The principal terms have the same
sign and modulus at least $cS(t)$; the noncenter assumption and
forcing exclude $S=0$.
If $u_0<0$, the condition $u>0$ gives $|u_0|\le Ct$.
If $z_0<0$, $z>0$ gives $|z_0|\le Cv\le Ct^r$.
In either case
\[
 |u_0z_0|\le CtS(t).
\]
Since $|Q'|\le Ct^{-1/8}$, its curvature contribution is bounded
by $C\epsilon|\sigma_s|t^{7/8}S(t)$.
Also $I\le C(|u_0|+|z_0|)$ and
$S(t)\ge t^{1/64}(|u_0|+|z_0|)$, whence
\[
 |\mathcal R_3'|\le CIt^{3/4}
                      \le Ct^{47/64}S(t).
\]
A common small $t_0$ makes both errors less than half the principal
margin. Therefore $M'$ cannot vanish.
\end{proof}

\paragraph{Ratios separated from one.}
Away from ratio one, a power comparison replaces the compensator.
We must also recover the signs of the original coefficients: the
curvature crossing alone does not determine them.
After reflection it is enough to fix $r\ge1+\delta$, $\delta>0$,
on a compact ratio interval. Put $I_a=|u_0|+|z_0|$.
The triangle-preserving coordinates give the unsheared equation
\[
 V_U=-rV/U-V^2R(U,V),\qquad
 R=\epsilon(u_0R_u+z_0R_z+\nu R_\nu),
\]
with bounded finite jets. Its exact reciprocal integral yields
\[
 v=t^re^w,\quad |D^jw|\le C\epsilon(I_a+|\nu|)t,\quad
 v'\asymp t^{r-1},\quad v''\ge c\,v'/t.
\]
Indeed $t^r\int_t^1s^{-r}\,\dd s\le Ct$, and its finitely many
Euler derivatives have the same bound. The differentiated integral
equations have small integral coefficients; iteration retains these
bounds. Fixed coefficient logarithms are absorbed with a portion of
the positive margin $\delta$. No near-resonant shear is used here.

Single-center division of the absolute closing difference gives
$|\nu|\le\delta_\epsilon I_a$ on every fibre with a closing zero.
Direct integral division of \eqref{rw:finite:three-M} then gives,
on the full positive $t$ interval of the auxiliary extension,
\begin{equation}\label{rw:finite:three-away-derivative}
 M'=u_0v'\mathcal A-z_0\mathcal B+\nu v\mathcal C,
\end{equation}
where
\[
 \mathcal A,\mathcal B>c>0,\quad
 |\mathcal A'|\le C,\quad
 |\mathcal B'|\le C(1+|u_0|/t),\quad
 |(v\mathcal C)'|\le Cv/t.
\]
These bounds follow by separating the two unmixed deficit terms
from their mixed connector term. On $u_0=z_0=0$, the numerator is
$tv\nu$ times a bounded coefficient, by division from the exact
triangle identity.
On $z_0=0$, divide $M$ by $v$ and then by $u_0$; its endpoint
coefficient is $(rh_--\epsilon)z_1>c$.
On $u_0=0$, divide by $t$ and then by $z_0$; its endpoint
coefficient is $-(h_+-\epsilon r)u_1<-c$.
The remaining mixed term has its constant
$\epsilon\kappa u_0z_0$ explicitly removed. Its first Euler
derivatives are $O(t)$ and its second ordinary derivative is
$O(1/t)$. The bounded connector derivatives and the preceding
bounds on $w$ prove these estimates term by term in
\eqref{rw:finite:three-M}; integral division in $u_0,z_0$
preserves them. Differentiating the two unmixed terms gives
\eqref{rw:finite:three-away-derivative}; the mixed second derivative
accounts for $|u_0|/t$ in $\mathcal B'$. Ratios of $v$ at points
of a coefficient integration segment are positive bounded units,
since $r$ is held fixed and the corresponding difference of $w$
is $O(\epsilon(I_a+|\nu|)t)$.

At a zero of $M'$, its two principal terms must balance. The
$\nu v$ term is smaller than their sum since
$v\le Ctv'$ and $v\le Ct$. Thus
\begin{equation}\label{rw:finite:three-away-balance}
 u_0z_0>0,\qquad |z_0|\asymp|u_0|v'
                      \asymp |u_0|t^{r-1}.
\end{equation}
At such a point,
$M''=u_0v''\mathcal A$ with relative error
$O(t+|u_0|+\delta_\epsilon t)$.
All zeros of $M'$ are simple and cross with sign $u_0$, so there
is at most one. This controls the auxiliary numerator. Actual returns
remain subject to the strict deficit conditions.

Suppose $M$ has two zeros in the comparison interval. The auxiliary
interval contains both zeros and extends to $t=0$. If $u_0,z_0$
were both negative, its only possible critical point would be a maximum, whereas
$M(0+)=\epsilon\kappa u_0z_0>0$. Such a function has at most one
positive zero. Hence
\begin{equation}\label{rw:finite:three-away-positive}
 u_0>0,\quad z_0>0,\quad
 z_0\le Cu_0t_*^{r-1}\le Cu_0t_0^\delta
\end{equation}
for a critical point $t_*$ between the two zeros. In particular,
after reducing $t_0$, $\xi_2=(u_0-z_0)/2>0$.

To recover $\tau_1$ with a uniform sign, work first at
$\epsilon=\nu=0$ on $z_0=0$. Let $q=q_l=Q_s$ be the common
incoming root. Its lower and upper equations are
\[
 \beta q^2+\bar a q=1,\qquad
 \bar m q^2-\bar c q=1.
\]
The outgoing intersections are $p_l=1/(\beta q)$ and
$P_s=1/(\bar m q)$. Therefore exactly
\begin{equation}\label{rw:finite:three-away-transverse}
 \tau_1=(\bar m-\beta)q,\qquad
 \delta_+=p_l-P_s=\frac{\bar m-\beta}{\beta\bar m q}.
\end{equation}
Since $u_0=\chi_+(\delta_+)$ with $\chi_+'>c$, these formulas
give $\tau_1=A_0u_0$ on this surface, where $A_0>c$ has bounded
jets. In particular the derivative in $\bar m-\beta$ is a
positive unit at the center, even when the actual ratio is used
as the remaining coordinate.

In the common inverse coordinates define, at positive scale,
$A_\epsilon=\int_0^1\partial_{u_0}
\tau_1(su_0,0,r,0;\epsilon)\,\dd s$.
The exact triangle identity makes the omitted constant zero.
Convergence of the lower strong intersections and the upper
low-layer jets to their zero-scale values gives
$|A_\epsilon-A_0|\le C\epsilon(1+|\log\epsilon|)$.
Single division in $z_0,\nu$ now gives
\[
 \tau_1=A_\epsilon u_0+B_\epsilon z_0+C_\epsilon\nu,\qquad
 A_\epsilon>c,\quad |B_\epsilon|+|C_\epsilon|\le C.
\]
Use \eqref{rw:finite:three-away-positive} and
$|\nu|\le\delta_\epsilon(u_0+z_0)$.
Choose $t_0$ with $C t_0^\delta$ small first, then the scale with
$C\delta_\epsilon$ small. The two errors are absorbed by
$A_\epsilon u_0/2$, and hence $\tau_1>0$. The lower and upper
intersection equations, together with the critical-point balance,
therefore determine the coefficient direction.
The same argument applied to $\bar m-\beta$ gives
$\bar m-\beta>0$: on the zero-scale surface in
\eqref{rw:finite:three-away-transverse} its quotient by $u_0$
is also a positive bounded unit, and it vanishes on the exact
positive-scale triangle axis. On the compact center part with
$r\ge1+\delta$, one has $\bar c=-\bar a<-c_\delta<0$;
the actual ratio coordinates preserve this margin on a smaller
neighborhood. Therefore the original primary generator has
$\sigma_1=\bar c(\bar m-\beta)<0$ on this nonresonant two-zero
branch. All $\nu$ terms in these divisions are $O(|\nu|)$ and
are absorbed by the stated forcing; no extra $\epsilon$ factor
is assigned to the $\nu$ dependence of the ratio.

The local calculation gives an upper bound of four closing zeros. It also identifies the
only branch on which the numerator can have two zeros. The signs on
that branch will be used at the already fixed regular seam in
Section~12; they do not yet determine the value of $\mathcal K$ at
that seam.

\begin{proposition}[Local count and the two-zero terminal branch]
\label{rw:finite:three-count}
On common small neighborhoods of every compact triangle-center
stratum, the closing equation has at most four isolated zeros on
\eqref{rw:finite:three-domain}, with multiplicity.
More precisely, on a noncenter fibre satisfying the local forcing
bound supplied by one closing zero,
\[
 Z(M)\le2,\qquad Z(\mathcal K)\le3,\qquad Z(F)\le4.
\]
If a terminal subinterval contains two zeros of $M$, counted with
multiplicity, orient the whole field by reflection if necessary so
that $r>1$. Then
\begin{equation}\label{rw:finite:three-terminal-signs}
 u_0>0,\qquad z_0>0,\qquad \xi_2>0,\qquad \tau_1>0,
 \qquad M(0+)>0,\qquad
 \mathcal K(t)\longrightarrow-\infty\quad(t\downarrow0).
\end{equation}
Near ratio one the same branch also has $T>0$, $W>0$ and
$\sigma_s<0$; away from one it has $\bar m>\beta$ and
$\sigma_1<0$. In particular $t=0$ is the end of the deficit
comparison interval in this branch. It need not be an end of the
positive-central return domain. If a seam $t_s$ in the comparison interval
is fixed away from the zeros and its tail contains exactly these two
zeros, then $M(t_s)>0$, $\mathcal K'_t(t_s)>0$, and
$\mathcal K'_p(t_s)<0$ in the outward coordinate
$p=P_s-\psi_+(t)$, $\psi_+'>0$.
These conclusions do not assume or provide a sign for
$\mathcal K(t_s)$.
\end{proposition}

\begin{proof}
Near one, \eqref{rw:finite:three-remainder} and
\eqref{rw:finite:three-column-norm} satisfy
Lemma~\ref{rw:finite:compensator-calculus}. Away from one, the
single crossing established above gives $Z(M)\le2$.
On the comparison interval, $M=tuz\,\mathcal K'$ with $tuz>0$,
and $F'=D_+'(u)G_+'(t)(1-e^{\mathcal K/\epsilon})$.
Two applications of Rolle's theorem with multiplicity therefore give
the displayed bounds. If there is no closing zero, the bound for $F$
is immediate. If forcing gives the exact center, there is no isolated
closing zero.

The away-from-one terminal signs were proved in
\eqref{rw:finite:three-away-positive} and
\eqref{rw:finite:three-away-transverse}. For the near-one signs,
first restrict to a smaller ratio interval so that its reflected
ratio remains in $|r-1|\le1/64$.
Fix a large $K_0$ and separate the cone
$|\xi_2|\le K_0|\kappa\xi_1|$.
On this cone, the exact positive column
\eqref{rw:finite:three-W-factor}, the small $\xi_2$ column,
and relative forcing give
\[
 \frac{W}{\kappa\xi_1}>c,\qquad
 |T/\xi_1|\le C|\kappa|,\qquad
 |\mathcal R_3'/\xi_1|\le C|\kappa|t^{1/2}.
\]
Thus the one-zero part of
Lemma~\ref{rw:finite:compensator-calculus} applies to $M/\xi_1$.
For $\kappa=0$ this cone is the reversible center. Neither case
permits the two-zero branch.

In the complementary cone choose $K_0$ large and then reduce the
coefficient box so that $T$ has the sign of $\xi_2$.
Two zeros of $M$ give a zero of $M'$ in the deficit interval, so
Lemma~\ref{rw:finite:three-critical-gaps} implies $u_0z_0>0$.
The case $\kappa=0$ would have $C_0=0$, contrary to
\eqref{rw:finite:compensator-signs}. For $\kappa>0$, $C_0>0$,
so that same lemma gives $W>0$ and $T>0$, hence $\xi_2>0$.
Equation~\eqref{rw:finite:three-directed-generators}, this cone and
forcing imply $\tau_1>0$ and $\sigma_s<0$.
The exact formula
\[
 W=\kappa u_0Q_*+2\epsilon u_0z_0\sigma_s r
\]
then excludes $u_0,z_0<0$: both terms would be negative.
Thus both deficits are positive.

Reflection exchanges the lower legs and replaces the actual upper
passage by its inverse. The new ratio is $1/r$, the new critical
function is $-\mathcal K$ composed with the inverse input, and its
derivative numerator is a nonzero signed unit times the old one.
It preserves zeros with multiplicity and the comparison interval.
Apply the preceding argument in that complete reflected coordinate
choice when $r<1$; no inequality is transported while keeping the
old entry coordinate fixed.

Positive endpoint deficits now place all sufficiently small $t>0$
in the comparison interval. Both complete central outputs must still
be positive for an actual return. At fixed coefficients, the lower
derivatives at $u_0,z_0$ are finite and nonzero, while
$v'\sim rt^{r-1}$.
This proves the endpoint limit for $\mathcal K$ in
\eqref{rw:finite:three-terminal-signs}.
Finally $M(0+)>0$ and exactly two zeros before the fixed seam,
counted with multiplicity, imply $M(t_s)>0$ by parity.
The positive multiplier and $dp/dt<0$ give the two derivative signs.
The remaining step is to determine $\mathcal K(t_s)$ from the
regular overlap. Proposition~\ref{rw:anchors:three-saddle} combines
that value with the derivative signs and endpoint parity to count
zeros on the whole union. The local estimate alone leaves the
four-zero closing bound.
\end{proof}
